% ============================================================
% PREGUNTAS POST-CLASE: TERMODINÁMICA
% Curso de Oriente - CCH Oriente
% Enero 2026
% ============================================================
% INSTRUCCIONES: Pegar este contenido en tu plantilla de Overleaf
% después de la definición de entornos y antes de \end{document}
% ============================================================
\renewcommand{\tema}{Preguntas Post-Clase - Termodinámica}

\twocolumn
\section*{PREGUNTAS POST-CLASE}

% ----------------------------------------------------------
% PREGUNTAS 1-5: Aplicación directa de conceptos
% ----------------------------------------------------------

\begin{preguntabox}
\subsection{Pregunta 1}
Un ingeniero mide la temperatura de los gases de escape de una turbina de gas y obtiene 527$^\circ$C. ¿Cuál es esa temperatura en kelvin?

\begin{enumerate}
    \item[A)] 254 K
    \item[B)] 527 K
    \item[C)] 800 K
    \item[D)] 1073 K
\end{enumerate}
\end{preguntabox}

\begin{preguntabox}
\subsection{Pregunta 2}
Una sonda espacial registra una temperatura de 77 K en la superficie de un satélite helado. ¿A cuántos grados Celsius corresponde esa lectura?

\begin{enumerate}
    \item[A)] $-$196$^\circ$C
    \item[B)] $-$350$^\circ$C
    \item[C)] 77$^\circ$C
    \item[D)] $-$77$^\circ$C
\end{enumerate}
\end{preguntabox}

\begin{preguntabox}
\subsection{Pregunta 3}
Un técnico en refrigeración necesita saber cuánto calor absorbe 500 g de agua cuando su temperatura pasa de 10$^\circ$C a 70$^\circ$C. ¿Cuántas kilocalorías se requieren? ($c_{\text{agua}} = 1$ cal/g$\cdot^\circ$C)

\begin{enumerate}
    \item[A)] 15 kcal
    \item[B)] 30 kcal
    \item[C)] 50 kcal
    \item[D)] 35 kcal
\end{enumerate}
\end{preguntabox}

\begin{preguntabox}
\subsection{Pregunta 4}
Un ciclista llena las llantas de su bicicleta en la madrugada cuando la temperatura es 7$^\circ$C y la presión es 4 atm. Al mediodía, la temperatura sube a 57$^\circ$C y el volumen permanece constante. ¿Cuál es la nueva presión? (Ley de Gay-Lussac)

\begin{enumerate}
    \item[A)] 4.4 atm
    \item[B)] 4.7 atm
    \item[C)] 5.0 atm
    \item[D)] 3.7 atm
\end{enumerate}
\end{preguntabox}

\begin{preguntabox}
\subsection{Pregunta 5}
Un submarino realiza una prueba de inmersión con un cilindro de aire a 10 atm y 60 L. Al soltar el gas a presión de 2 atm a temperatura constante (proceso isotérmico), ¿qué volumen ocupa el aire?

\begin{enumerate}
    \item[A)] 120 L
    \item[B)] 200 L
    \item[C)] 300 L
    \item[D)] 600 L
\end{enumerate}
\end{preguntabox}

% ----------------------------------------------------------
% PREGUNTAS 6-10: Combinación de varios conceptos
% ----------------------------------------------------------

\begin{preguntabox}
\subsection{Pregunta 6}
Un bloque de aluminio de 200 g se calienta 50$^\circ$C en un horno de fundición. ¿Cuántas calorías absorbe? ($c_{\text{Al}} = 0.215$ cal/g$\cdot^\circ$C)

\begin{enumerate}
    \item[A)] 1 075 cal
    \item[B)] 2 150 cal
    \item[C)] 4 300 cal
    \item[D)] 860 cal
\end{enumerate}
\end{preguntabox}

\begin{preguntabox}
\subsection{Pregunta 7}
En el laboratorio de un hospital, un gas ideal ocupa 6 L a 300 K y 2 atm. Se calienta a 400 K y se comprime hasta 3 atm. ¿Cuál es el nuevo volumen?

\begin{enumerate}
    \item[A)] 4 L
    \item[B)] 6 L
    \item[C)] 8 L
    \item[D)] 3 L
\end{enumerate}
\end{preguntabox}

\begin{preguntabox}
\subsection{Pregunta 8}
El compresor de una planta industrial inyecta aire en un recipiente rígido. Al inicio el gas está a 27$^\circ$C y 3 atm. Después de operar, la temperatura llega a 127$^\circ$C. ¿Cuánto aumentó la presión respecto al valor inicial?

\begin{enumerate}
    \item[A)] 1 atm
    \item[B)] 3 atm
    \item[C)] 4 atm
    \item[D)] 0.5 atm
\end{enumerate}
\end{preguntabox}

\begin{preguntabox}
\subsection{Pregunta 9}
Un gas ideal en un émbolo absorbe 1 200 J de calor y su energía interna aumenta en 500 J. ¿Cuánto trabajo realizó el gas al expandirse? (Primera ley)

\begin{enumerate}
    \item[A)] 500 J
    \item[B)] 700 J
    \item[C)] 1 200 J
    \item[D)] 1 700 J
\end{enumerate}
\end{preguntabox}

\begin{preguntabox}
\subsection{Pregunta 10}
Una pieza de hierro de 1 000 g sale de un torno a 173$^\circ$C y se sumerge en agua hasta enfriarse a 23$^\circ$C. ¿Cuántas kilocalorías cedió la pieza de hierro? ($c_{\text{Fe}} = 0.113$ cal/g$\cdot^\circ$C)

\begin{enumerate}
    \item[A)] 11.3 kcal
    \item[B)] 16.95 kcal
    \item[C)] 22.6 kcal
    \item[D)] 5.65 kcal
\end{enumerate}
\end{preguntabox}

% ----------------------------------------------------------
% PREGUNTAS 11-15: Problemas desafiantes con múltiples pasos
% ----------------------------------------------------------

\begin{preguntabox}
\subsection{Pregunta 11}
Un globo aerostático contiene 600 L de aire a 27$^\circ$C y 1 atm. Al encender el quemador, la presión se mantiene constante y la temperatura sube a 177$^\circ$C. ¿Cuántos litros adicionales ocupa el aire respecto al volumen inicial?

\begin{enumerate}
    \item[A)] 150 L
    \item[B)] 200 L
    \item[C)] 300 L
    \item[D)] 900 L
\end{enumerate}
\end{preguntabox}

\begin{preguntabox}
\subsection{Pregunta 12}
Un panel solar suministra 400 kcal de calor a 2 kg de agua a 25$^\circ$C. ¿A qué temperatura final llega el agua? ($c_{\text{agua}} = 1$ cal/g$\cdot^\circ$C)

\begin{enumerate}
    \item[A)] 175$^\circ$C
    \item[B)] 200$^\circ$C
    \item[C)] 225$^\circ$C
    \item[D)] 425$^\circ$C
\end{enumerate}
\end{preguntabox}

\begin{preguntabox}
\subsection{Pregunta 13}
El motor de un tractor absorbe 5 000 J de calor de la combustión y cede 3 000 J al ambiente. Según la primera ley de la termodinámica, ¿cuánto trabajo útil realiza el motor?

\begin{enumerate}
    \item[A)] 8 000 J
    \item[B)] 3 000 J
    \item[C)] 2 000 J
    \item[D)] 5 000 J
\end{enumerate}
\end{preguntabox}

\begin{preguntabox}
\subsection{Pregunta 14}
Una lata de refresco de aluminio de 150 g se enfría de 30$^\circ$C a $-$10$^\circ$C en una nevera portátil. ¿Cuántas calorías cedió la lata al refrigerante? ($c_{\text{Al}} = 0.215$ cal/g$\cdot^\circ$C)

\begin{enumerate}
    \item[A)] 645 cal
    \item[B)] 1 290 cal
    \item[C)] 967.5 cal
    \item[D)] 430 cal
\end{enumerate}
\end{preguntabox}

\begin{preguntabox}
\subsection{Pregunta 15}
Un aerogenerador almacena aire comprimido a 5 atm y 27$^\circ$C en un depósito de 400 L. Durante la noche, la temperatura baja a $-$73$^\circ$C. Si el volumen no cambia, ¿cuál es la presión del aire almacenado?

\begin{enumerate}
    \item[A)] 5 atm
    \item[B)] 2.5 atm
    \item[C)] 3 atm
    \item[D)] 4 atm
\end{enumerate}
\end{preguntabox}

% ----------------------------------------------------------
% SECCIÓN DE RESPUESTAS
% ----------------------------------------------------------

\newpage
\onecolumn
\section*{RESPUESTAS Y RETROALIMENTACIÓN}

\begin{respuestabox}
\subsection{Pregunta 1}
\textcolor{verdecorrecto}{\textbf{Respuesta correcta: C) 800 K}}

\textbf{Justificación:}\\
\textbf{Fórmula utilizada:} Conversión Celsius a Kelvin (Lección: Sección 1 --- Escalas Termométricas)
$$T(\text{K}) = T(^\circ\text{C}) + 273$$

\textbf{Sustitución:}
$$T(\text{K}) = 527\,^\circ\text{C} + 273 = \mathbf{800\,\text{K}}$$

\textbf{Respuestas incorrectas:}
\begin{itemize}
    \item \textbf{A) 254 K:} se restó 273 en lugar de sumar; la fórmula correcta suma la constante al valor en Celsius.
    \item \textbf{B) 527 K:} se copió el número de Celsius sin convertir; ignorar la constante de conversión es un error frecuente.
    \item \textbf{D) 1073 K:} se duplicó la temperatura en Celsius antes de sumar, operación sin fundamento.
\end{itemize}
\end{respuestabox}

\begin{respuestabox}
\subsection{Pregunta 2}
\textcolor{verdecorrecto}{\textbf{Respuesta correcta: A) $-$196$^\circ$C}}

\textbf{Justificación:}\\
\textbf{Fórmula utilizada:} Conversión Kelvin a Celsius (Lección: Sección 1 --- Escalas Termométricas)
$$T(^\circ\text{C}) = T(\text{K}) - 273$$

\textbf{Sustitución:}
$$T(^\circ\text{C}) = 77\,\text{K} - 273 = \mathbf{-196\,^\circ\text{C}}$$

\textbf{Respuestas incorrectas:}
\begin{itemize}
    \item \textbf{B) $-$350$^\circ$C:} se multiplicó por un factor incorrecto en lugar de restar 273.
    \item \textbf{C) 77$^\circ$C:} se confundió la escala; 77 es el valor en kelvin, no en Celsius. Son escalas distintas.
    \item \textbf{D) $-$77$^\circ$C:} se cambió el signo del valor original sin aplicar la fórmula de conversión.
\end{itemize}
\end{respuestabox}

\begin{respuestabox}
\subsection{Pregunta 3}
\textcolor{verdecorrecto}{\textbf{Respuesta correcta: B) 30 kcal}}

\textbf{Justificación:}\\
\textbf{Fórmula utilizada:} Calor transferido (Lección: Sección 4 --- Calor Específico y Capacidad Calorífica)
$$Q = m \cdot c \cdot \Delta T$$

\textbf{Sustitución:}
$$Q = 500\,\text{g} \times 1\,\frac{\text{cal}}{\text{g}\cdot^\circ\text{C}} \times (70 - 10)\,^\circ\text{C}$$
$$Q = 500\,\text{g} \times 1\,\frac{\text{cal}}{\text{g}\cdot^\circ\text{C}} \times 60\,^\circ\text{C} = 30{,}000\,\text{cal} = \mathbf{30\,\text{kcal}}$$

\textbf{Respuestas incorrectas:}
\begin{itemize}
    \item \textbf{A) 15 kcal:} se dividió el resultado entre 2; posible confusión con la temperatura final dividida a la mitad.
    \item \textbf{C) 50 kcal:} se usó $\Delta T = 70\,^\circ$C (temperatura final) en lugar de la diferencia $T_f - T_i = 60\,^\circ$C.
    \item \textbf{D) 35 kcal:} se multiplicó por $\Delta T = 70$ sin restarle la temperatura inicial.
\end{itemize}
\end{respuestabox}

\begin{respuestabox}
\subsection{Pregunta 4}
\textcolor{verdecorrecto}{\textbf{Respuesta correcta: C) 5.0 atm}}

\textbf{Justificación:}\\
\textbf{Fórmula utilizada:} Ley de Gay-Lussac (Lección: Sección 6 --- Ley de Gay-Lussac)
$$\frac{P_1}{T_1} = \frac{P_2}{T_2} \qquad (V = \text{constante})$$

Primero se convierten las temperaturas a kelvin:
$$T_1 = 7 + 273 = 280\,\text{K} \qquad T_2 = 57 + 273 = 330\,\text{K}$$

\textbf{Sustitución:}
$$P_2 = P_1 \cdot \frac{T_2}{T_1} = 4\,\text{atm} \times \frac{330\,\text{K}}{280\,\text{K}} = \mathbf{5.0\,\text{atm}}$$

\textbf{Respuestas incorrectas:}
\begin{itemize}
    \item \textbf{A) 4.4 atm:} se usaron las temperaturas en Celsius (57/7) sin convertir a Kelvin.
    \item \textbf{B) 4.7 atm:} error aritmético al sustituir; posiblemente se usó $T_2 = 330$ K con $T_1 = 280$ K pero se dividió incorrectamente.
    \item \textbf{D) 3.7 atm:} se invirtió la razón de temperaturas: se calculó $P_2 = P_1 \times T_1/T_2$ en lugar de $T_2/T_1$.
\end{itemize}
\end{respuestabox}

\begin{respuestabox}
\subsection{Pregunta 5}
\textcolor{verdecorrecto}{\textbf{Respuesta correcta: C) 300 L}}

\textbf{Justificación:}\\
\textbf{Fórmula utilizada:} Ley de Boyle (Lección: Sección 6 --- Ley de Boyle)
$$P_1 V_1 = P_2 V_2 \qquad (T = \text{constante})$$

\textbf{Despeje:}
$$V_2 = \frac{P_1 V_1}{P_2}$$

\textbf{Sustitución:}
$$V_2 = \frac{10\,\text{atm} \times 60\,\text{L}}{2\,\text{atm}} = \mathbf{300\,\text{L}}$$

\textbf{Respuestas incorrectas:}
\begin{itemize}
    \item \textbf{A) 120 L:} se sumó presión y volumen en lugar de multiplicar; operación sin sentido dimensional.
    \item \textbf{B) 200 L:} se dividió solo el volumen entre el cociente de presiones sin el factor correcto; error de manipulación algebraica.
    \item \textbf{D) 600 L:} se multiplicó el volumen por la presión inicial sin dividir entre la presión final.
\end{itemize}
\end{respuestabox}

\begin{respuestabox}
\subsection{Pregunta 6}
\textcolor{verdecorrecto}{\textbf{Respuesta correcta: B) 2 150 cal}}

\textbf{Justificación:}\\
\textbf{Fórmula utilizada:} Calor transferido (Lección: Sección 4 --- Calor Específico y Capacidad Calorífica)
$$Q = m \cdot c \cdot \Delta T$$

\textbf{Sustitución:}
$$Q = 200\,\text{g} \times 0.215\,\frac{\text{cal}}{\text{g}\cdot^\circ\text{C}} \times 50\,^\circ\text{C} = \mathbf{2{,}150\,\text{cal}}$$

\textbf{Respuestas incorrectas:}
\begin{itemize}
    \item \textbf{A) 1 075 cal:} se dividió el resultado entre 2; equivale a haber tomado solo la mitad de la masa o la mitad del calor específico.
    \item \textbf{C) 4 300 cal:} se multiplicó el resultado correcto por 2; posible confusión con el calor específico del agua (1 cal/g$\cdot^\circ$C) en lugar del de aluminio.
    \item \textbf{D) 860 cal:} se omitió $\Delta T = 50\,^\circ$C y se multiplicó solo $m \times c$ sin la variación de temperatura.
\end{itemize}
\end{respuestabox}

\begin{respuestabox}
\subsection{Pregunta 7}
\textcolor{verdecorrecto}{\textbf{Respuesta correcta: A) 4 L}}

\textbf{Justificación:}\\
\textbf{Fórmula utilizada:} Ley combinada de gases ideales (Lección: Sección 6 --- Ley General de los Gases Ideales)
$$\frac{P_1 V_1}{T_1} = \frac{P_2 V_2}{T_2}$$

\textbf{Despeje:}
$$V_2 = V_1 \cdot \frac{P_1}{P_2} \cdot \frac{T_2}{T_1}$$

\textbf{Sustitución:}
$$V_2 = 6\,\text{L} \times \frac{2\,\text{atm}}{3\,\text{atm}} \times \frac{400\,\text{K}}{300\,\text{K}} = 6\,\text{L} \times \frac{2}{3} \times \frac{4}{3} = 6\,\text{L} \times \frac{8}{9} \approx \mathbf{4\,\text{L}}$$

\textbf{Respuestas incorrectas:}
\begin{itemize}
    \item \textbf{B) 6 L:} se supuso que el volumen no cambia, ignorando que tanto la presión como la temperatura varían.
    \item \textbf{C) 8 L:} se aplicó solo el factor de temperatura ($T_2/T_1$) sin considerar el aumento de presión que comprime el gas.
    \item \textbf{D) 3 L:} se aplicó solo el factor de presión ($P_1/P_2$) sin considerar el efecto del calentamiento.
\end{itemize}
\end{respuestabox}

\begin{respuestabox}
\subsection{Pregunta 8}
\textcolor{verdecorrecto}{\textbf{Respuesta correcta: A) 1 atm}}

\textbf{Justificación:}\\
\textbf{Fórmula utilizada:} Ley de Gay-Lussac (Lección: Sección 6 --- Ley de Gay-Lussac)
$$\frac{P_1}{T_1} = \frac{P_2}{T_2} \qquad (V = \text{constante})$$

Conversión de temperaturas:
$$T_1 = 27 + 273 = 300\,\text{K} \qquad T_2 = 127 + 273 = 400\,\text{K}$$

\textbf{Sustitución:}
$$P_2 = 3\,\text{atm} \times \frac{400\,\text{K}}{300\,\text{K}} = 4\,\text{atm}$$

El \textbf{aumento} de presión es:
$$\Delta P = P_2 - P_1 = 4\,\text{atm} - 3\,\text{atm} = \mathbf{1\,\text{atm}}$$

\textbf{Respuestas incorrectas:}
\begin{itemize}
    \item \textbf{B) 3 atm:} se reportó la presión inicial como el aumento, sin calcular $\Delta P$.
    \item \textbf{C) 4 atm:} se reportó la presión final $P_2$ como si fuera el aumento; la pregunta pide la diferencia.
    \item \textbf{D) 0.5 atm:} error aritmético; posiblemente se calculó $(T_2 - T_1)/T_1$ con los valores en Celsius sin convertir.
\end{itemize}
\end{respuestabox}

\begin{respuestabox}
\subsection{Pregunta 9}
\textcolor{verdecorrecto}{\textbf{Respuesta correcta: B) 700 J}}

\textbf{Justificación:}\\
\textbf{Fórmula utilizada:} Primera ley de la termodinámica (Lección: Sección 5 --- Primera Ley)
$$\Delta U = Q - W$$

\textbf{Despeje del trabajo:}
$$W = Q - \Delta U$$

\textbf{Sustitución:}
$$W = 1{,}200\,\text{J} - 500\,\text{J} = \mathbf{700\,\text{J}}$$

\textbf{Respuestas incorrectas:}
\begin{itemize}
    \item \textbf{A) 500 J:} se reportó $\Delta U$ como el trabajo; se confundió la variación de energía interna con el trabajo realizado.
    \item \textbf{C) 1 200 J:} se supuso que todo el calor absorbido se convirtió en trabajo, lo que violaría la primera ley si $\Delta U \neq 0$.
    \item \textbf{D) 1 700 J:} se sumó $Q + \Delta U$ en lugar de restar; signo incorrecto en la fórmula.
\end{itemize}
\end{respuestabox}

\begin{respuestabox}
\subsection{Pregunta 10}
\textcolor{verdecorrecto}{\textbf{Respuesta correcta: B) 16.95 kcal}}

\textbf{Justificación:}\\
\textbf{Fórmula utilizada:} Calor transferido (Lección: Sección 4 --- Calor Específico y Capacidad Calorífica)
$$Q = m \cdot c \cdot \Delta T$$

El hierro baja de temperatura, así que cede calor ($\Delta T = 150\,^\circ$C):
$$\Delta T = 173\,^\circ\text{C} - 23\,^\circ\text{C} = 150\,^\circ\text{C}$$

\textbf{Sustitución:}
$$Q = 1{,}000\,\text{g} \times 0.113\,\frac{\text{cal}}{\text{g}\cdot^\circ\text{C}} \times 150\,^\circ\text{C} = 16{,}950\,\text{cal} = \mathbf{16.95\,\text{kcal}}$$

\textbf{Respuestas incorrectas:}
\begin{itemize}
    \item \textbf{A) 11.3 kcal:} se usó $\Delta T = 100\,^\circ$C en lugar de los 150$^\circ$C correctos; posible confusión con la diferencia hasta 73$^\circ$C.
    \item \textbf{C) 22.6 kcal:} se multiplicó el resultado por 2; equivale a haber usado $\Delta T = 300\,^\circ$C.
    \item \textbf{D) 5.65 kcal:} se dividió el resultado entre 3; error aritmético sin justificación física.
\end{itemize}
\end{respuestabox}

\begin{respuestabox}
\subsection{Pregunta 11}
\textcolor{verdecorrecto}{\textbf{Respuesta correcta: C) 300 L}}

\textbf{Justificación:}\\
\textbf{Fórmula utilizada:} Ley de Charles (Lección: Sección 6 --- Ley de Charles)
$$\frac{V_1}{T_1} = \frac{V_2}{T_2} \qquad (P = \text{constante})$$

Conversión de temperaturas:
$$T_1 = 27 + 273 = 300\,\text{K} \qquad T_2 = 177 + 273 = 450\,\text{K}$$

\textbf{Sustitución:}
$$V_2 = V_1 \cdot \frac{T_2}{T_1} = 600\,\text{L} \times \frac{450\,\text{K}}{300\,\text{K}} = 900\,\text{L}$$

El volumen adicional respecto al inicial es:
$$\Delta V = V_2 - V_1 = 900\,\text{L} - 600\,\text{L} = \mathbf{300\,\text{L}}$$

\textbf{Respuestas incorrectas:}
\begin{itemize}
    \item \textbf{A) 150 L:} se calculó $\Delta V$ usando temperaturas en Celsius (177/27) sin convertir a Kelvin.
    \item \textbf{B) 200 L:} error de proporción; posiblemente se dividió 600 L entre 3 sin fundamento.
    \item \textbf{D) 900 L:} se reportó el volumen final $V_2$ como el volumen adicional; la pregunta pide el incremento $\Delta V$, no el valor total.
\end{itemize}
\end{respuestabox}

\begin{respuestabox}
\subsection{Pregunta 12}
\textcolor{verdecorrecto}{\textbf{Respuesta correcta: C) 225$^\circ$C}}

\textbf{Justificación:}\\
\textbf{Fórmula utilizada:} Calor transferido, despeje de temperatura final (Lección: Sección 4 --- Calor Específico y Capacidad Calorífica)
$$Q = m \cdot c \cdot \Delta T \implies T_f = T_i + \frac{Q}{m \cdot c}$$

Datos: $Q = 400\,\text{kcal} = 400{,}000\,\text{cal}$, $m = 2{,}000\,\text{g}$, $c = 1\,\text{cal/(g}\cdot^\circ\text{C)}$, $T_i = 25\,^\circ\text{C}$.

\textbf{Sustitución:}
$$\Delta T = \frac{400{,}000\,\text{cal}}{2{,}000\,\text{g} \times 1\,\text{cal/(g}\cdot^\circ\text{C)}} = 200\,^\circ\text{C}$$
$$T_f = 25\,^\circ\text{C} + 200\,^\circ\text{C} = \mathbf{225\,^\circ\text{C}}$$

\textbf{Respuestas incorrectas:}
\begin{itemize}
    \item \textbf{A) 175$^\circ$C:} se restó $\Delta T$ de $T_i$ en lugar de sumarlo.
    \item \textbf{B) 200$^\circ$C:} se reportó $\Delta T$ como la temperatura final, olvidando sumar la temperatura inicial de 25$^\circ$C.
    \item \textbf{D) 425$^\circ$C:} se sumó $T_i + \Delta T$ usando la masa en kg (2) sin convertir a gramos; el denominador quedó 100 veces más chico.
\end{itemize}
\end{respuestabox}

\begin{respuestabox}
\subsection{Pregunta 13}
\textcolor{verdecorrecto}{\textbf{Respuesta correcta: C) 2 000 J}}

\textbf{Justificación:}\\
\textbf{Fórmula utilizada:} Primera ley de la termodinámica (Lección: Sección 5 --- Primera Ley)
$$\Delta U = Q - W$$

En un ciclo de motor, la energía interna no cambia ($\Delta U = 0$) porque el motor regresa a su estado inicial al completar el ciclo. Por conservación de energía:
$$Q_H = W + Q_C \implies W = Q_H - Q_C$$

\textbf{Sustitución:}
$$W = 5{,}000\,\text{J} - 3{,}000\,\text{J} = \mathbf{2{,}000\,\text{J}}$$

\textbf{Respuestas incorrectas:}
\begin{itemize}
    \item \textbf{A) 8 000 J:} se sumó $Q_H + Q_C$ en lugar de restar; la energía no se crea, se transforma.
    \item \textbf{B) 3 000 J:} se reportó el calor cedido al ambiente ($Q_C$) como si fuera el trabajo útil.
    \item \textbf{D) 5 000 J:} se reportó el calor absorbido total ($Q_H$) como trabajo, ignorando las pérdidas; eso equivaldría a eficiencia del 100\%, lo que viola la segunda ley.
\end{itemize}
\end{respuestabox}

\begin{respuestabox}
\subsection{Pregunta 14}
\textcolor{verdecorrecto}{\textbf{Respuesta correcta: B) 1 290 cal}}

\textbf{Justificación:}\\
\textbf{Fórmula utilizada:} Calor transferido (Lección: Sección 4 --- Calor Específico y Capacidad Calorífica)
$$Q = m \cdot c \cdot |\Delta T|$$

La lata cede calor al enfriarse, así que $\Delta T = 30 - (-10) = 40\,^\circ$C.

\textbf{Sustitución:}
$$Q = 150\,\text{g} \times 0.215\,\frac{\text{cal}}{\text{g}\cdot^\circ\text{C}} \times 40\,^\circ\text{C} = \mathbf{1{,}290\,\text{cal}}$$

\textbf{Respuestas incorrectas:}
\begin{itemize}
    \item \textbf{A) 645 cal:} se usó $\Delta T = 30\,^\circ$C (solo temperatura inicial) en lugar de la diferencia total de 40$^\circ$C.
    \item \textbf{C) 967.5 cal:} se usó $\Delta T = 30\,^\circ$C y se multiplicó por $150 \times 0.215$; la temperatura final negativa se ignoró.
    \item \textbf{D) 430 cal:} se omitió la masa y se calculó $c \times \Delta T$ sin multiplicar por los 150 g.
\end{itemize}
\end{respuestabox}

\begin{respuestabox}
\subsection{Pregunta 15}
\textcolor{verdecorrecto}{\textbf{Respuesta correcta: B) 2.5 atm}}

\textbf{Justificación:}\\
\textbf{Fórmula utilizada:} Ley de Gay-Lussac (Lección: Sección 6 --- Ley de Gay-Lussac)
$$\frac{P_1}{T_1} = \frac{P_2}{T_2} \qquad (V = \text{constante})$$

Conversión de temperaturas:
$$T_1 = 27 + 273 = 300\,\text{K} \qquad T_2 = -73 + 273 = 200\,\text{K}$$

\textbf{Sustitución:}
$$P_2 = P_1 \cdot \frac{T_2}{T_1} = 5\,\text{atm} \times \frac{200\,\text{K}}{300\,\text{K}} = \mathbf{2.5\,\text{atm}}$$

\textbf{Respuestas incorrectas:}
\begin{itemize}
    \item \textbf{A) 5 atm:} se asumió que la presión no cambia al bajar la temperatura; eso solo ocurre si el gas puede expandirse, lo que no es posible en un depósito rígido.
    \item \textbf{C) 3 atm:} error aritmético; posiblemente se usó una razón de temperaturas incorrecta (200/300 redondeado a 0.6 y aplicado sin la multiplicación correcta).
    \item \textbf{D) 4 atm:} se invirtió la razón de temperaturas: se calculó $5 \times (300/200)$ que daría 7.5, o bien se usó una razón parcial errónea.
\end{itemize}
\end{respuestabox}